\chapter{Future Work}
\label{chap:future-work}

The results in this thesis show that task-graph information can be used for
locality-aware placement and guarded DVFS control inside CUDASTF. At the same
time, the current implementation is deliberately conservative and the
evaluation is limited to a single-node benchmark suite. This chapter outlines
the most important directions for extending the work.

\section{Broader Hardware and Workload Evaluation}
\label{sec:future-broader-evaluation}

The first priority is broader evaluation. The current results are from one
eight-GPU node and five reported benchmarks. Future experiments should repeat the
placement and DVFS evaluation on different GPU generations, different
interconnect topologies, and different problem sizes. This would test whether
the locality improvements remain stable when peer-copy bandwidth, copy-engine
contention, memory capacity, and computation time change.

The benchmark suite should also be expanded. The current workloads include
strong locality winners and neutral controls, but a larger suite would make it
easier to identify the boundary between useful residual placement and
HEFT-like fallback behavior. In particular, applications with irregular data
reuse, phase changes, and mixed compute/communication behavior would be useful
stress tests for both the placement gate and the DVFS guardrails.

\section{More Accurate Slack and Criticality Models}
\label{sec:future-slack-criticality}

The current scheduler uses a local criticality proxy based on task cost,
dependency count, and input-byte volume. This proxy is inexpensive and suitable
for online use, but it is not a full dynamic critical-path analysis. A natural
extension is to add a lightweight DAG-aware slack estimator that updates as the
runtime observes task completions and queue states.

Such a model could improve both layers of the scheduler. For placement, more
accurate slack estimates could allow the residual gate to distinguish between
safe non-HEFT deviations and deviations that are likely to delay the critical
path. For DVFS, better slack information could identify windows where lower
frequency is safe because the extra runtime is hidden by dependency structure
or communication. The challenge is to gain this information without adding
enough overhead to erase the benefit.

\section{Topology-Aware Locality Modelling}
\label{sec:future-topology-locality}

The current locality model treats copy saving as a central scheduling signal,
but the copy-cost estimate is intentionally lightweight. Future work should
make this model topology aware. On modern multi-GPU systems, the cost of
moving data depends on whether GPUs are connected by NVLink, PCIe, or another
interconnect, and on whether copy engines are already busy.

A topology-aware model could assign different costs to different source--target
GPU pairs and could account for asymmetric bandwidth or contention. It could
also distinguish copy bytes that are likely to overlap with computation from
copy bytes that are likely to appear on the critical path. This would make the
locality certificate more precise while preserving the HEFT-relative structure
of the scheduler.

\section{Richer DVFS Objectives}
\label{sec:future-dvfs-objectives}

The WinDVFS layer in this thesis focuses on guarded energy reduction over the
benchmark execution window with bounded latency overhead. Future work could
extend the committer to optimize different objectives, such as EDP or a
user-defined energy budget. The same proposal--commit structure could support
these objectives by changing the reward or eligibility rule used at window
boundaries.

Another direction is to learn the DVFS commit policy online. The current
committer is rule based: it uses proposal shares, criticality mass, backlog,
slack budget, and slowdown guardrails. A learned committer could adapt these
thresholds across workloads while keeping the existing guardrails as hard
safety constraints. This would preserve the conservative behavior needed for
runtime use while reducing manual parameter sensitivity.

The DVFS action space could also be expanded. The current design primarily
controls graphics frequency and treats memory frequency conservatively. Some
workloads may benefit from coordinated graphics- and memory-clock decisions,
especially when memory bandwidth rather than core throughput dominates
runtime. Adding memory-clock control would require stronger safeguards because
incorrect memory-frequency decisions can increase latency substantially.

\section{Transition Placement and Overhead Reduction}
\label{sec:future-overhead}

The proposal--commit design amortizes frequency decisions over windows, but it
does not fully solve transition placement. Future work should measure and model
where transitions become visible in the CUDASTF execution timeline. If the
runtime can schedule transitions during communication, synchronization, or
other slack periods, the visible latency cost of DVFS could be reduced.

Instrumentation overhead should also be reduced. The static-replay diagnostic
in this thesis suggests that online Bandit decision overhead is small at
benchmark scale, but the replay experiment is still only an approximation of a
zero-overhead scheduler. A production-oriented implementation should separate
low-overhead counters used during normal execution from detailed trace logging
used only for debugging and thesis-style attribution. This would keep the
scheduler explainable without making detailed logs part of the critical path.

\section{Reproducibility and Deployment}
\label{sec:future-reproducibility}

The final direction is reproducibility and deployment. Future versions should
standardize the experiment records produced by each run, including benchmark
summaries, scheduler configuration, device metadata, and power-sampling
summaries. A stable record format would make it easier to compare results
across hardware and to regenerate thesis figures without relying on ad hoc log
parsing.

For deployment, the scheduler should expose a small set of user-facing modes:
for example, placement-only, conservative DVFS, and diagnostic logging. These
modes would let users choose between maximum stability, energy savings, and
explainability depending on the application. The long-term goal is a CUDASTF
scheduler that offers locality and energy-efficiency modes by default while
remaining transparent enough for users to understand when and why it deviates
from HEFT.
